% Part: computability
% Chapter: computability-theory
% Section: total

\documentclass[../../../include/open-logic-section]{subfiles}

\begin{document}

\olfileid{cmp}{thy}{tot}
\olsection{全体性是不可判定的}

下面再举一个运用 $s$-$m$-$n$ 定理证明不可计算性的例子。
令 $\fn{Tot}$ 为全可计算函数的指标集，即
\[
\fn{Tot} = \Setabs{x}{\text{对每个 $y$，$\cfind{x}(y)\fdefined$}}.
\]

\begin{prop}
\ollabel{prop:total}
$\fn{Tot}$ 不是可计算的。
\end{prop}

\begin{proof}
要证明 $\fn{Tot}$ 不可计算，只需证明 $K$ 可归约到它。
定义 $h(x,y)$ 为
\[
h(x,y) \simeq
\begin{cases}
0 & \text{若 $x \in K$} \\
\fundefined & \text{否则}
\end{cases}
\]
注意 $h(x,y)$ 完全不依赖于 $y$。
不难看出 $h$ 是部分可计算的：对于输入 $x, y$，计算~$h$ 的过程是先模拟函数~$\cfind{x}$ 在输入~$x$ 上的计算；
如果该计算停机，则 $h(x,y)$ 输出 $0$ 并停机。
所以 $h(x,y)$ 就是 $\Zero(\umin{s}{T(x,x,s)})$，其中 $\Zero$ 是常值零函数。

利用 $s$-$m$-$n$ 定理，存在一个原始递归函数~$k(x)$，使得对每个 $x$ 和~$y$，
\[
\cfind{k(x)}(y) =
\begin{cases}
0 & \text{若 $x \in K$} \\
\fundefined & \text{否则}
\end{cases}
\]
因此，若 $x \in K$ 则 $\cfind{k(x)}$ 是全函数，否则无定义。
于是，$k$ 是从 $K$ 到~$\fn{Tot}$ 的归约。
\end{proof}

\begin{digress}
实际上，$\fn{Tot}$ 甚至不是可计算枚举的——它的复杂度位于“算术分层”的更上方。
但我们在此不讨论这一更强的结论。
\end{digress}

\end{document}